\scanpage{26}
\mspage{24}

3. \uline{Essai de construction du foncteur} $\Omega$

Il n'y a sans doute guère de chances de prouver quelque chose sur le foncteur

\[
 \Omega:MW^{-1}\longrightarrow MW^{-1},
\]

sans \unsure{fournir|donner} au moins une \uline{construction} plus ou moins explicite de celui-ci. Je vais essayer, \deleted{\unsure{en}} en m'inspirant de Quillen, de dégager des conditions générales sur $M,W$ (plus générales que celles de Quillen) qui assurent l'existence de $\Omega$, et permettent de le construire. L'idée géométrique remonte à Serre-Cartan. \deleted{c'est} Il s'agit, pour un objet $a$ de $M$, de construire un objet \deleted{\unsure{pour}} $\mathrm{Ch}(\deleted{\unsure{a}}S)$ \uline{dans $M$}, jouant le rôle d'un \deleted{\unsure{tel}} ``espace de chemins'', et qui doit \deleted{\unsure{satisfaire à la condition}} \inserted{\unsure{s'inclure|s'insérer} dans un diagramme de} factorisation de \deleted{\unsure{\(\Delta\)}} $\mathrm{diag}_S:S\longrightarrow S\times S$ (on suppose que dans $M$ les \deleted{produits} \inserted{\unsure{lims|lims.}} finies, \inserted{donc les produits finis,} existent)

\diagram{\begin{tikzcd}[ampersand replacement=\&,row sep=small,column sep=small]
 S\arrow[d,"i_S"]\arrow[dd,bend right=22,"{\mathrm{diag}_S}"']\&\\
 \mathrm{Ch}(S)\arrow[d,"{\unsure{p|P}_S}"]\&\\
 S\times S\&e\arrow[l]
\end{tikzcd}}

La propriété essentielle, c'est que $p$ soit une ``\uline{fibration}'' (dans un sens convenable, à dégager).
% archive sheet number: 25
